\section{Related Work}
\label{sec:related_work}

\subsection{Handwritten Text Recognition}
\label{sec:rw_htr}

Two architecture families dominate modern offline HTR. The first is the CNN--BiLSTM--CTC family, canonically represented by the ``Best Practices'' recipe of Retsinas~\etal~\cite{retsinas2022htrbestpractices}: aspect-ratio-preserving preprocessing, a convolutional stem that collapses the image height to a one-dimensional token sequence, a bidirectional recurrent head, and Connectionist Temporal Classification (CTC)~\cite{graves2006ctc} to align the variable-length transcription with a fixed-length output grid. This family remains competitive on limited training data because the CNN's built-in inductive biases match the physical structure of ink strokes. A survey of the broader field, including HMM-era systems and the transition to end-to-end deep models, is given by Memon~\etal~\cite{memon2020htrsurvey}.

The second family replaces the CNN encoder with a Vision Transformer. HTR-VT~\cite{li2024htrvt} keeps a CNN stem for the two-dimensional-to-one-dimensional reduction and delegates the sequence encoding to a plain ViT encoder with span masking, again supervised by CTC. TrOCR~\cite{li2023trocr} instead uses a pretrained ViT encoder--decoder trained on hundreds of millions of real and synthetic handwriting lines. Both models are used as external comparators in this report; neither integrates register tokens or analyses attention quality under a register-count sweep.

\subsection{Vision Transformers and Register Tokens}
\label{sec:rw_registers}

The Transformer~\cite{vaswani2017attention} and its Vision Transformer variant~\cite{dosovitskiy2021vit} replace convolutional feature extraction with multi-head self-attention over a flattened patch sequence. In a classification ViT a special CLS token is prepended to the patch sequence, self-attention aggregates image content into it across layers, and a linear head produces the class posterior.

Darcet~\etal~\cite{darcet2024registers} identified a failure mode of the resulting attention maps. In well-trained supervised and self-supervised ViTs, a small number of patch tokens develop anomalously high norms and disproportionate attention mass, acting as \emph{attention sinks} that produce visible artifacts and hurt any downstream task that reads the attention. Prepending a small number of learnable \emph{register tokens} to the patch sequence moves the sinks onto the registers, leaving patch attention clean. The mechanism is architecturally minimal --- it adds only $R \cdot D$ parameters for register count $R$ and embedding width $D$ --- and it leaves classification accuracy essentially unchanged.

Two structural differences separate the classification-ViT setting from the CTC HTR setting used in this report. First, a CTC ViT has no dedicated CLS query: the queries that produce the CTC-decoded output are the patch tokens themselves. Second, the CTC loss is agnostic to attention structure and does not directly penalize a map that ``looks'' bad if the greedy decoding is correct. Both differences suggest that the register-token effect on attention may not transfer cleanly, and this is the specific question the report evaluates.

\subsection{Attention-Based Interpretability and Downstream Writer Analysis}
\label{sec:rw_interp}

The interest in clean HTR attention maps is grounded in two downstream applications. The first is attention-based character localization for synthetic handwriting generation, developed by Nikolaidou~\etal\ in \emph{Beyond Memorization}~\cite{nikolaidou2024beyondmem} on top of the WordStylist pipeline~\cite{nikolaidou2023wordstylist}. Their Figure~5 shows per-character attention maps localized on individual glyphs, produced by the cross-attention layers of a diffusion U-Net in which explicit text-token queries attend to image-region keys. The visualization style is the one this report reproduces in Section~\ref{sec:res_bm}. Two architectural caveats are relevant: the reference maps come from cross-attention rather than self-attention, and their input is a near-square word crop rather than an aspect-elongated HTR line, so an exact pixel-level reproduction is not possible --- what is reproducible is the visual grammar of a per-character attention overlay that shifts spatially with the recognized character.

The second application is writer identification via character-indexed feature aggregation. Souibgui~\etal~\cite{souibgui2022vlac} introduce Vectors of Locally Aggregated Characters (VLAC), a page-level writer descriptor that aggregates local descriptors extracted at each recognized character position, weighted by the recognizer's per-character attention. The method is interpretable at the character level and is the direct downstream justification for the ``clean attention'' target of this project. A full VLAC evaluation is out of scope here and is sketched as future work in Section~\ref{sec:conclusion}.

Positioning-wise, this report combines three elements that no prior work combines: a CTC-based sequence-labelling loss, a register-token intervention, and a per-character attention-quality analysis of the resulting maps. The methodological substrate is Retsinas~\etal~\cite{retsinas2022htrbestpractices}; the encoder follows the design used in HTR-VT~\cite{li2024htrvt}; the register mechanism follows Darcet~\etal~\cite{darcet2024registers}; and the visualization target is Nikolaidou~\etal~\cite{nikolaidou2024beyondmem} Fig.~5, with the honest caveat that our attention type differs from theirs.
